\chapter{Coherent multi-qubit sensing and the Heisenberg limit}
\label{ch:coherent-multiqubit}

The four mechanism chapters of Part~VI
(\cref{ch:qp-detection,ch:qubit-antenna,ch:phonon-qubit-piezo,ch:microwave-absorption})
all read out a \emph{single} qubit, or at most a small number of
independent qubits, after a dark-matter--induced excitation
event. They take the qubit as a binary energy meter --- did a
quasiparticle break a Cooper pair? did an itinerant photon land
in the cavity? --- and integrate the count rate against the
dark-count rate. There is, however, a second mode of operation
that is uniquely available to the qubit substrate: the
dark-matter field $\phi(t)$ at the chip imprints a small but
\emph{coherent} phase on every qubit it touches, and that phase
can in principle be extracted with a precision that improves not
as $1/\sqrt{N}$ in the number of probes (the Ramsey-like
\emph{standard quantum limit}, SQL) but as $1/N$ (the
\emph{Heisenberg limit}, HL). Realising the Heisenberg
scaling requires entangled probes and is fragile against
decoherence; the connection to quantum error correction
(\cref{ch:quantum-information-processing}) is exactly that QEC
is the tool that protects the Heisenberg-limited phase-sensing
protocol from the same Markovian noise that, in qubit
information processing, would corrupt logical states. This
closing chapter of Part~VI explains the SQL and HL, derives the
$1/N$ scaling for an entangled-qubit phase estimator, states the
Zhou--Zhang--Preskill--Jiang theorem on when QEC restores
the Heisenberg limit in the presence of noise, and points to
the recent dark-matter experiments that put these ideas into
hardware.

% ============================================================
\section{Phase imprinted by a dark-matter field}
\label{sec:dm-phase}

\begin{phenomenon}[\textbf{Dark matter as a coherent phase signal}]
A bosonic dark-matter field at a single fixed point in the
laboratory --- a dark photon $A'_\mu$, an axion field $a$, or
any weakly coupled wave-like dark-matter candidate ---
oscillates at angular frequency $\omega_\chi=m_\chi c^2/\hbar$
with coherence time $\tau_\chi\sim 1/(m_\chi v^2/2)$ (set by the
$10^{-3}c$ velocity dispersion of the local halo). Through its
coupling to the qubit Hamiltonian (e.g.\ a kinetically mixed
electric field on the antenna of \cref{ch:qubit-antenna}, a
piezoelectrically driven strain in
\cref{ch:phonon-qubit-piezo}, or a coherent intracavity photon
field in \cref{ch:microwave-absorption}), the dark-matter
field shifts the qubit transition frequency by a small,
oscillating amount $\delta\omega_q(t)\propto\phi(t)$. Over a
sensing time $T<\tau_\chi$, this accumulates a coherent
phase $\varphi$ that can be read out by Ramsey
interferometry. The question of \emph{how precisely
$\varphi$ can be estimated} is the question of how well the
underlying coupling $\varepsilon$ (or $g_{a\gamma\gamma}$) can
be bounded.
\end{phenomenon}

\paragraph{From mechanism to phase.}
The four mechanism chapters all gave the dark-matter field a
seat at the qubit Hamiltonian: in
\cref{ch:qubit-antenna,ch:microwave-absorption} via an
electromagnetic coupling, in \cref{ch:phonon-qubit-piezo} via
strain. Schematically, a single qubit feels
\begin{equation}
\label{eq:dm-phase-hamiltonian}
\op H(t) \;=\; -\frac{\hbar\,\omega_q}{2}\,\hat\sigma_z
\;+\; \hbar\,g\,\phi(t)\,\hat\sigma_z,
\end{equation}
where $g$ is the relevant coupling (units of frequency per unit
field amplitude). In a frame rotating at $\omega_q$, the second
term integrates over a sensing window $[0,T]$ to a coherent
phase
\begin{equation}
\label{eq:dm-phase-integral}
\varphi(T) \;=\; g \int_0^T\! \phi(t)\,dt \;\approx\;
\frac{g\,\phi_0}{\omega_\chi}\,
\bigl(\sin\omega_\chi T\bigr)
\;\xrightarrow[\omega_\chi T\ll 1]{}\;
g\,\phi_0\,T,
\end{equation}
i.e.\ in the long-coherence limit ($T\ll\tau_\chi$ and
$\omega_\chi T\lesssim 1$) the phase grows linearly in $T$ and
linearly in the dark-matter field amplitude $\phi_0$. Estimating
$\varphi$ to precision $\sigma_\varphi$ thus translates directly
into a bound on $\phi_0$, hence on the underlying coupling
$\varepsilon$ or $g_{a\gamma\gamma}$.

% ============================================================
\section{Independent probes: the standard quantum limit}
\label{sec:sql}

\begin{statement}[\textbf{Standard quantum limit (SQL)}]
\label{stmt:sql}
For $N$ independent two-level probes, each prepared in the
equal superposition $\ket{+}=(\ket{0}+\ket{1})/\sqrt 2$,
exposed for time $T$ to the Hamiltonian
\cref{eq:dm-phase-hamiltonian}, and read out projectively in
the $\hat\sigma_x$ basis, the optimal estimator $\hat\varphi$
of the accumulated phase $\varphi$ saturates the
Cram\'{e}r--Rao bound
\begin{equation}
\label{eq:sql-bound}
(\Delta\varphi)_{\rm SQL} \;=\; \frac{1}{\sqrt{N}},
\qquad
(\Delta\phi_0)_{\rm SQL} \;=\; \frac{1}{g\,T\,\sqrt{N}}\,.
\end{equation}
The same $1/\sqrt{N}$ scaling applies to a single probe queried
$N$ times sequentially, provided each query is independent and
the noise is Markovian.
\end{statement}

\paragraph{Sketch derivation.}
Each independent qubit, after evolution under
\cref{eq:dm-phase-hamiltonian} for time $T$, is in the state
$\ket{\psi(T)}=(\ket{0}+e^{-i\varphi}\ket{1})/\sqrt 2$.
Measurement in $\hat\sigma_x$ gives outcome $+1$ with
probability $p(\varphi)=\cos^2(\varphi/2)$. The Fisher information
per qubit is~\cite{degen2017}
\begin{equation}
F_1(\varphi) \;=\; \frac{[\partial_\varphi p]^2}{p\,(1-p)} \;=\; 1,
\end{equation}
and additivity of Fisher information across independent probes
gives $F_N=N$, so the Cram\'{e}r--Rao bound reads
$(\Delta\varphi)^2 \geq 1/N$. This is the same $1/\sqrt N$
that appears in classical statistics: the central-limit
inheritance of measurement scatter over uncorrelated
samples~\cite{caves1981,wineland1992}.

\paragraph{Why this matters for dark matter.}
With a single transmon ($N=1$) coupled to a dark photon at
$g\,\phi_0\sim 10^{-15}\,\omega_q$ for a $T=100\,\mu$s
Ramsey window, the accumulated phase is $\varphi\sim 10^{-12}$,
far below the per-shot $\Delta\varphi=1$. Repeating the
experiment $N\sim 10^{18}$ times --- a year of integration
at MHz repetition --- recovers the dark-matter phase only at
SQL scaling. Any tool that improves the scaling exponent from
$1/\sqrt N$ to $1/N$ saves a factor of $\sqrt{N}$ in
integration time at fixed sensitivity, which can be
many orders of magnitude.

% ============================================================
\section{Entangled probes: the Heisenberg limit}
\label{sec:hl}

\begin{phenomenon}[\textbf{The Heisenberg limit}]
If the $N$ qubits are not independent but are prepared in a
particular entangled state --- the Greenberger--Horne--Zeilinger
(GHZ) state ---
\begin{equation}
\label{eq:ghz}
\ket{\rm GHZ} \;=\; \frac{1}{\sqrt 2}\bigl(\ket{0\cdots 0}
+ \ket{1\cdots 1}\bigr),
\end{equation}
then the same single-particle phase $\varphi$ stamped on each
qubit shifts the global state phase by $N\varphi$, not
$\varphi$. The interferometer arms have effectively become
$N$ times longer, the fringe is $N$ times finer, and the
phase-estimation precision improves as $1/N$. This is the
``Heisenberg-limited'' (HL) scaling first identified for
photonic interferometers~\cite{caves1981,giovannetti2004,giovannetti2006}
and for atomic-clock spectroscopy~\cite{wineland1992,huelga1997}.
\end{phenomenon}

\begin{statement}[\textbf{Heisenberg limit (HL)}]
\label{stmt:hl}
For $N$ qubits prepared in $\ket{\rm GHZ}$ and exposed to the
Hamiltonian \cref{eq:dm-phase-hamiltonian}, the
Cram\'{e}r--Rao bound on the phase estimate is
\begin{equation}
\label{eq:hl-bound}
(\Delta\varphi)_{\rm HL} \;=\; \frac{1}{N},
\qquad
(\Delta\phi_0)_{\rm HL} \;=\; \frac{1}{g\,T\,N}\,.
\end{equation}
The improvement over the SQL is a factor of $\sqrt{N}$ in
phase resolution, equivalently a factor of $N$ in integration
time at fixed sensitivity. The bound is saturated by an optimal
estimator of the form $\hat\varphi=\arg\langle X_{\rm GHZ}\rangle/N$,
where $X_{\rm GHZ}=\bigotimes_i\hat\sigma_x^{(i)}$ is the
parity operator that converts the global $N\varphi$ phase into
a single fringe.
\end{statement}

\paragraph{Derivation.}
Apply $\op U(T)=\exp[-i T \op H/\hbar]$ from
\cref{eq:dm-phase-hamiltonian} (dropping the static
$\omega_q$ term, which is removed by the rotating frame) to
the GHZ state \cref{eq:ghz}. Each qubit picks up
$e^{-i\varphi/2}$ on $\ket 1$ and $e^{+i\varphi/2}$ on
$\ket 0$, so
\begin{equation}
\op U(T)\ket{\rm GHZ} \;=\; \frac{1}{\sqrt 2}
\bigl(e^{+iN\varphi/2}\ket{0\cdots 0}
+ e^{-iN\varphi/2}\ket{1\cdots 1}\bigr),
\end{equation}
giving a relative phase $N\varphi$ between the two
many-body branches. The expectation of the parity operator
$X_{\rm GHZ}$ in this state is $\cos(N\varphi)$, with
variance $\sin^2(N\varphi)$, and the Cram\'{e}r--Rao bound
delivers \cref{eq:hl-bound}. The factor of $N$ in the phase
\emph{is} the entanglement: $N$ qubits acting as a single
$N$-fold-amplified phase pickup~\cite{giovannetti2011,pezze2018}.

\paragraph{Squeezing as an alternative route to sub-SQL scaling.}
GHZ states are not the only quantum-enhanced phase sensors.
Spin-squeezed states of $N$ qubits (or $N$ spins of an atomic
ensemble) reduce the noise on one quadrature of the collective
spin at the cost of increased noise on the conjugate
quadrature; the result is sub-SQL scaling, parameterised by
the Wineland squeezing parameter $\xi^2<1$~\cite{wineland1992}.
Quadrature-squeezed light injected into a haloscope plays the
same role for an electromagnetic phase: HAYSTAC's two-mode
squeezed-vacuum operation~\cite{malnou2019,haystac2021} uses
this exact mechanism to accelerate axion searches.

\begin{example}[\textbf{Heisenberg-limited dark-photon phase pickup}]
\label{ex:hl-darkphoton}
Consider $N=10$ transmon qubits, each with $T_2=100\,\mu$s,
coupled to a dark photon with $g\phi_0/2\pi=1$\,Hz. In a single
Ramsey shot, each qubit accumulates $\varphi=2\pi\,(g\phi_0\,T_2)
\sim 6\times 10^{-4}$\,rad. With independent probes
(\cref{stmt:sql}) the phase resolution per shot is
$\Delta\varphi_{\rm SQL}=1/\sqrt{10}\approx 0.32$\,rad, so the
signal sits $\sim 500$ shots below the noise floor. With the
qubits in $\ket{\rm GHZ}_{10}$ and read out by parity
(\cref{stmt:hl}) the phase resolution is
$\Delta\varphi_{\rm HL}=1/10=0.1$\,rad, a factor of
$\sqrt{10}\approx 3.2$ better. Maintained for the dark-matter
coherence time $\tau_\chi\sim 1$\,ms (at $\omega_\chi/2\pi
\sim 5$\,GHz), this saves a factor of 10 in integration time at
fixed sensitivity --- modest at $N=10$, but the saving scales
linearly in $N$.
\end{example}

\paragraph{Transition.}
The arithmetic of \cref{ex:hl-darkphoton} is, however, an
\emph{ideal-case} arithmetic: it neglects the very dephasing,
relaxation, and readout noise that real superconducting qubits
suffer from on the same $T_2\sim 100\,\mu$s timescale. The next
section confronts this directly: GHZ states are exquisitely
fragile, and Markovian noise demolishes the Heisenberg scaling
back to SQL. The way out --- and the bridge to Part~V on quantum
information processing --- is quantum error correction.

% ============================================================
\section{Decoherence, the no-go theorem, and the QEC restoration}
\label{sec:qec-metrology}

\begin{phenomenon}[\textbf{Decoherence reverts the HL to SQL}]
A single qubit subject to Markovian dephasing at rate
$\Gamma_\phi$ has a coherence time $T_2=1/\Gamma_\phi$ over
which a phase can be accumulated before the off-diagonal
density-matrix element decays. A GHZ state of $N$ qubits, by
contrast, decoheres $N$ times faster: its off-diagonal element
between $\ket{0\cdots 0}$ and $\ket{1\cdots 1}$ decays as
$e^{-N\Gamma_\phi t}$, because every qubit contributes
independently. The signal grows by $N$, the noise grows by
$\sqrt{N}\cdot\sqrt{N}=N$, and the net advantage cancels:
asymptotically the HL collapses back onto the
SQL~\cite{huelga1997,escher2011,demkowicz2012}.
\end{phenomenon}

\begin{statement}[\textbf{Zhou--Zhang--Preskill--Jiang HNLS theorem}]
\label{stmt:hnls}
Let the noisy quantum probe evolve under a Markovian master
equation with signal Hamiltonian $\op G$ (the generator whose
parameter $\varphi=g\phi_0 T$ we wish to estimate) and noise
described by a set of Lindblad operators $\{\op L_k\}$ acting
on a logical sub-system supported by a quantum
error-correcting code. Assume that fast,
noiseless ancillas and noiseless quantum operations on the
ancilla--probe joint system are available. Then the Heisenberg
scaling $(\Delta\varphi)\sim 1/N$ is achievable by a QEC sensing
protocol if and only if
\begin{equation}
\label{eq:hnls}
\op G \;\notin\; \op{LS}\bigl(\{\op L_k,\op L_k^\dagger,
\op L_k^\dagger\op L_k,\id\}\bigr),
\end{equation}
i.e.\ the signal generator $\op G$ does \emph{not} lie in
the Lindblad span of the noise operators (the ``HNLS''
condition: \emph{Hamiltonian-Not-in-Lindblad-Span}). When
HNLS holds, an explicit QEC code exists that simultaneously
preserves the signal Hamiltonian and projects out the noise;
when HNLS fails, no protocol --- regardless of how clever ---
beats the SQL asymptotically~\cite{zhou2018}.
\end{statement}

\paragraph{Reading the HNLS condition.}
The condition \cref{eq:hnls} has a clean physical
interpretation. The Lindblad span is the Hilbert-space subspace
of operators that can be assembled out of the noise operators
$\op L_k$ and the identity. Any signal generator that lives
\emph{inside} that span looks, to the system, like a noise
operator: the same QEC code that catches and reverses the
$\op L_k$ errors will also catch and reverse the signal,
erasing it. By contrast, a signal generator that is
\emph{orthogonal} to the noise span carries information that
the code cannot mistake for noise, and the QEC cycle preserves
the signal Hamiltonian while sterilising the noise. For a
qubit subject to pure dephasing along $\hat\sigma_z$
($\op L_1=\sqrt{\Gamma_\phi/2}\,\hat\sigma_z$) and a signal that
also commutes with $\hat\sigma_z$ (the dark-matter case of
\cref{eq:dm-phase-hamiltonian}!), HNLS is \emph{not}
satisfied: the signal lies in the Lindblad span and no QEC
code can restore the HL --- the dephasing-imposed SQL is
fundamental in this orientation. This is the noise/signal
geometry that explains why the dark-matter phase coupling on a
qubit \emph{should not, naively, expect} a Heisenberg
advantage. HL is recovered if the signal generator is rotated
relative to the noise --- e.g.\ by encoding the qubit in a
non-trivial multi-mode code (the cat code, the GKP code, or a
distance-2 stabiliser code) so that the dephasing axis no
longer coincides with the signal axis~\cite{kessler2014,arrad2014,dur2014,zhou2018}.

\begin{statement}[\textbf{HL restoration via QEC --- protocol sketch}]
\label{stmt:qec-sensing}
The construction of \cite{zhou2018} runs a QEC cycle on
the $(N+\rm{ancilla})$-qubit system at rate $\gg \Gamma_\phi$.
At each cycle:
\begin{enumerate}
  \item the syndrome of the noise operators $\op L_k$ is
        measured (via stabiliser measurements);
  \item the syndrome is flipped on the ancilla and used to
        Pauli-correct the data qubits;
  \item the residual evolution is dominated by the
        signal Hamiltonian $\op G$, which the code preserves
        because $\op G\notin\op{LS}$.
\end{enumerate}
After $T/T_{\rm QEC}\gg 1$ cycles, the encoded logical phase
has accumulated $\varphi_L \sim N\,g\,\phi_0\,T$ with residual
noise that does \emph{not} grow as $N$, recovering
the HL scaling \cref{eq:hl-bound}.
\end{statement}

\paragraph{Connecting to Part~V (quantum information processing).}
The conceptual bridge to QIP is direct: the QEC codes that
protect logical qubits during a quantum computation are the
\emph{same} codes that protect a logical phase sensor during a
quantum-enhanced measurement, with the syndrome-extraction
cycle playing the same role in both contexts. The
\emph{distance} of the code (minimum weight of an undetectable
error) controls how many simultaneous noise events can be
tolerated before the encoded signal is lost; the
\emph{transversality} structure of the code controls which
generators $\op G$ are protected. In the language of Part~V,
the chapter on quantum error correction states the threshold
theorem: arbitrary fault tolerance is achievable for
computation if and only if physical error rates are below a
constant. The Zhou--Zhang--Preskill--Jiang theorem is the
\emph{metrological} counterpart: arbitrary HL scaling is
achievable for sensing if and only if the signal lies outside
the noise's Lindblad span~\cite{zhou2018,layden2019}. Both are
necessary-and-sufficient conditions on the geometry of noise
relative to the operation one wants to perform.

% ============================================================
\section{Outlook: dark matter at the Heisenberg limit}
\label{sec:dm-hl-outlook}

\begin{application}[\textbf{Quantum-enhanced dark-matter searches}]
\label{app:dm-hl}
Several experimental directions translate the SQL/HL
distinction into hardware:
\begin{itemize}
  \item \textbf{Squeezed-vacuum haloscopes.} HAYSTAC operates
        a microwave Josephson parametric amplifier in
        two-mode-squeezed-vacuum mode, demonstrating a factor of
        $\sim 2$ scan-rate improvement over the SQL operating
        point at $\sim 5$\,GHz~\cite{malnou2019,haystac2021}.
        This is a single-mode realisation of the spin-squeezing
        idea applied to the haloscope's signal mode rather than
        to a multi-qubit register.
  \item \textbf{Entangled sensor networks.} A network of
        $N$ qubits, each at a different physical location and
        each tuned to the same dark-matter Compton frequency,
        can in principle be entangled across the network and
        read out via a global parity measurement. The
        Heisenberg-limited sensitivity of such a distributed
        haloscope to ultralight bosonic dark matter has been
        analysed in detail~\cite{brady2022},
        with operational demonstrations beginning to appear in
        2024~\cite{liu2024,connolly2024}.
  \item \textbf{Single-mode error-corrected sensors.}
        Bosonic codes (cat codes, binomial codes, GKP codes)
        encode a logical phase in a high-photon-number cavity
        mode and have been used for single-mode QEC of an
        oscillator~\cite{ofek2016,sivak2023}. Adapting them to
        run as a Heisenberg-limited dark-matter phase sensor
        is an active programme of theoretical and experimental
        work~\cite{layden2019}.
\end{itemize}
The unifying picture is that a dark-matter haloscope, whether
microwave-cavity or qubit-based, is at heart a phase
interferometer: the target field $\phi(t)$ paints a tiny
oscillating phase onto the apparatus, and the sensitivity to
$\varepsilon$ or $g_{a\gamma\gamma}$ is set by how
precisely that phase can be estimated. Quantum-enhanced
metrology --- squeezing, entanglement, error correction ---
attacks that bottleneck directly, and so far only at the
percent--$\sqrt N$ level. The conceptual ceiling is the
Heisenberg limit; the engineering ceiling is the HNLS-protected
QEC sensor; both remain open frontiers.
\end{application}

\paragraph{Concluding outlook.}
Across Part~VI we have followed the dark-matter signal from
its halo physics (\cref{ch:dm-search}), through the
four cQED hardware mechanisms by which it can deposit energy
or phase on a chip
(\cref{ch:qp-detection,ch:qubit-antenna,ch:phonon-qubit-piezo,ch:microwave-absorption}),
to the metrological architecture by which the deposited signal
is converted into a coupling bound. The key conceptual
takeaway of this closing chapter is that the qubit substrate
buys two distinct kinds of dark-matter sensitivity: an
\emph{event}-counting sensitivity, set by the qubit dark-count
rate, and a \emph{phase}-sensing sensitivity, set in the
quantum-enhanced limit by the Heisenberg-scaling theorems and
the HNLS condition. The former drives the recoil and absorption
chapters; the latter is the natural home for the squeezed
haloscopes, entangled sensor networks, and error-corrected
phase sensors that close the gap between the SQL and the
fundamental quantum limit. A fault-tolerant
dark-matter sensor --- one whose precision is limited only by
the dark-matter coherence time itself, not by the qubit's
$T_2$ --- is the natural endpoint of this programme, and
it is the same fault-tolerance theorem that powers a
fault-tolerant quantum computer
(\cref{ch:quantum-information-processing}).
